% !TEX program = xelatex
% 每个环境的用法示例；改动 yupin-notes.cls 后重新编译本文件检查。
\documentclass{yupin-notes}
\coursetitle{高等概率论}
\coursekind{课堂笔记}
\coursesubtitle{Durrett《Probability: Theory and Examples》第 1–5 章}
\courseyear{2026–2027 学年秋季学期}
\courseauthor{王汉超}
\courseinstitution{山东大学}
\homeworkheader{山东大学 2026 年秋季学期，高等概率论，作业 #1}

\begin{document}
\makecoursetitle
\tableofcontents

\lecture{示例：概率空间，$\sigma$-代数，测度的连续性}
\label{lec:demo}

一段导入正文。术语第一次出现时加黑并标注英文：\term{$\sigma$-代数}（$\sigma$-algebra）。

\topic{$\sigma$-代数}

\begin{summary}
方框导语：本节要做什么，为什么。
\end{summary}

\begin{definition}[$\sigma$-代数]\label{def:demo:sigma}
设 $\Omega$ 为非空集合，$\cF\subseteq 2^\Omega$。若 (1) $\Omega\in\cF$；(2) $A\in\cF\Rightarrow A^c\in\cF$；(3) $A_1,A_2,\dots\in\cF\Rightarrow\bigcup_n A_n\in\cF$，则称 $\cF$ 为 $\Omega$ 上的 \marginterm{$\sigma$-代数}。
\end{definition}

\begin{example}
$\{\varnothing,\Omega\}$ 与 $2^\Omega$ 都是 $\sigma$-代数。
\end{example}

\begin{remark}
注记不编号。
\end{remark}

\begin{theorem}[测度的连续性]\label{thm:demo:cont}
若 $A_n\uparrow A$，则 $P(A_n)\uparrow P(A)$。于是 $X_n\toas X$ 蕴含 $X_n\toP X$，且 $E\abs{X}=\int\abs{X}\dd P$，$\Var(X)=E X^2-(EX)^2$，$\one_A$ 为示性函数。
\end{theorem}
\begin{proof}
令 $B_1=A_1$，$B_n=A_n\setminus A_{n-1}$……由定义~\ref{def:demo:sigma} 知 $B_n\in\cF$。
\end{proof}

\begin{theorem*}[Carathéodory 扩张定理]
不编号的定理。
\end{theorem*}

\begin{practice}
正文中的小练习。
\end{practice}

\homework{测度的连续性与 Borel–Cantelli 引理}

\exercisepart{基本习题}
\begin{exercises}{A}{课堂内容的补充}
  \item\label{ex:demo:A1} 证明定理~\ref{thm:demo:cont} 的下降版本。
  \item 第二题。
\end{exercises}

\begin{probleme}{B}{Borel–Cantelli 引理}
引导语。
\begin{questions}
  \item 第一问。
  \item 第二问。
  \begin{questions}\item 子问。\end{questions}
\end{questions}
过渡文字。
\begin{questions}[resume]
  \item 第三问，可用 \exref{ex:demo:A1} 的结论。
\end{questions}
\end{probleme}

\thinkingpart
\begin{probleme}*[问题]{C}{Vitali 不可测集}
难题。
\begin{questions}\item 第一问。\end{questions}
\end{probleme}

\exercisepart{课外知识：Kolmogorov 0-1 律}
\begin{probleme}{D}{}
\begin{questions}\item 第一问。\end{questions}
\end{probleme}
\end{document}
